\section{Conclusion}
\label{sec:conclusion}

We asked whether a still-safe prefix of a \qwen response reveals that a clear safety violation is coming within the next 64--512 tokens. On a family-split benchmark of natural generations, it partly does. A linear probe on a 16-token activation history reaches \auroc 0.79 at 64 tokens and 0.90 at 512 tokens, and it flags 56\% of unsafe trajectories at least 128 tokens before onset at a 5\% benign false-alarm rate. An off-the-shelf Activation Oracle does poorly zero-shot. Training on future-event labels closes most of that gap, but an explicit forecast representation adds nothing over a matched direct reader. The key takeaway is that future-event supervision, not forecast architecture, drives the gains, and that a cheap activation-history probe is the baseline any verbalizer-based forecaster must beat.

Future work should scale the number of positive families with more rollouts per harmful prompt, and calibrate labels with repeated rollouts. Windowed or recurrent encoders could exploit the history signal in both the direct and forecast readers. The most direct next test is an LLM-decoder AO fine-tuned on these future-event labels, with and without the forecast bottleneck. Model organisms with controlled delays, compared against a token-position control, would show whether the history signal reflects plans or position.
